On ne considere dans cette section que le probleme de l'evolution
de la distribution des electrons dans l'espace des configurations.
On omet donc ici la dynamique des flux en $p$ et $\xi.$ C'est l'approche
classique prise en compte dans tous les codes Fokker-Planck 3D actuels.
Les valeurs des fonctions et des flux sont evaluees a la meme position
$p$ et $\xi$ dans l'espace des impulsions, soit a la demi-grille
en impulsion $p_{i+1/2}$ et en diffusion angulaire $\xi_{j+1/2}$.
Elle presente l'avantage de la simplicite numerique, le transport
dans l'espace des configurations ne couplant que des electrons ayant
les memes valeurs en $p_{\Vert}$ et $p_{\bot}.$ L'energie des electrons
proportionnelle a $p^{2}$ est donc bien conservee, mais le moment
magnetique, qui est en principe un invariant du mouvement, n'est en
revanche pas conserve. Compte tenu des hypotheses, on omettra de specifier
les indices $i+1/2$ et $j+1/2$. L'equation de transport de base
est donnee par \[
\frac{\partial\left(rp^{2}f_{0}\right)}{\partial t}+p^{2}\frac{\partial\left(rS_{r}\right)}{\partial r}+p^{2}\frac{\partial S_{\theta}}{\partial\theta}=0\]
 avec \[
S_{r}=-D_{rr}\frac{\partial f_{0}}{\partial r}-\frac{1}{r}D_{r\theta}\frac{\partial f_{0}}{\partial\theta}+F_{r}f_{0}\]
 \[
S_{\theta}=-D_{\theta r}\frac{\partial f_{0}}{\partial r}-\frac{1}{r}D_{\theta\theta}\frac{\partial f_{0}}{\partial\theta}+F_{\theta}f_{0}\]


Sous forme discrete celle-ci devient \begin{eqnarray*}
r_{l+1/2}p_{i+1/2}^{2}\frac{f_{0,l+1/2,m+1/2}^{(k+1)}-f_{0,l+1/2,m+1/2}^{(k)}}{\Delta t} & = & -p_{i+1/2}^{2}\frac{r_{l+1}\left.S_{r}\right|_{l+1,m+1/2}^{(k+1)}-r_{l}\left.S_{r}\right|_{l,m+1/2}^{(k+1)}}{\Delta r_{l+1/2}}\\
 &  & -p_{i+1/2}^{2}\frac{\left.S_{\theta}\right|_{l+1/2,m+1}^{(k+1)}-\left.S_{\theta}\right|_{l+1/2,m}^{(k+1)}}{\Delta\theta_{m+1/2}}\end{eqnarray*}
 avec \begin{eqnarray*}
\left.S_{r}\right|_{l+1,m+1/2}^{(k+1)} & = & -D_{rr,l+1,m+1/2}\left.\frac{\partial f_{0}}{\partial r}\right|_{l+1,m+1/2}^{(k+1)}\\
 &  & -D_{r\theta,l+1,m+1/2}\frac{1}{r_{l+1}}\left.\frac{\partial f_{0}}{\partial\theta}\right|_{l+1,m+1/2}^{(k+1)}+F_{r,l+1,m+1/2}f_{0,l+1,m+1/2}^{(k+1)}\\
\left.S_{r}\right|_{l,m+1/2}^{(k+1)} & = & -D_{rr,l,m+1/2}\left.\frac{\partial f_{0}}{\partial r}\right|_{l,m+1/2}^{(k+1)}\\
 &  & -D_{r\theta,l,m+1/2}\frac{1}{r_{l}}\left.\frac{\partial f_{0}}{\partial\theta}\right|_{l,m+1/2}^{(k+1)}+F_{r,l,m+1/2}f_{0,l,m+1/2}^{(k+1)}\\
\left.S_{\theta}\right|_{l+1/2,m+1}^{(k+1)} & = & -D_{\theta r,l+1/2,m+1}\left.\frac{\partial f_{0}}{\partial r}\right|_{l+1/2,m+1}^{(k+1)}\\
 &  & -D_{\theta\theta,l+1/2,m+1}\frac{1}{r_{l+1/2}}\left.\frac{\partial f_{0}}{\partial\theta}\right|_{l+1/2m+1}^{(k+1)}+F_{\theta,l+1/2,m+1}f_{0,l+1/2,m+1}^{(k+1)}\\
\left.S_{\theta}\right|_{l+1/2,m}^{(k+1)} & = & -D_{\theta r,l+1/2,m}\left.\frac{\partial f_{0}}{\partial r}\right|_{l+1/2,m}^{(k+1)}\\
 &  & -D_{\theta\theta,l+1/2,m}\frac{1}{r_{l+1/2}}\left.\frac{\partial f_{0}}{\partial\theta}\right|_{l+1/2m}^{(k+1)}+F_{\theta,l+1/2,m}f_{0,l+1/2,m}^{(k+1)}\end{eqnarray*}
 Les expressions discretes des derivees partielles prennent la forme
suivante, \[
\left\{ \begin{array}{l}
\left.\frac{\partial f_{0}}{\partial r}\right|_{l+1,m+1/2}^{(k+1)}=\frac{f_{0,l+3/2,m+1/2}^{(k+1)}-f_{0,l+1/2,m+1/2}^{(k+1)}}{\Delta r_{l+1}}\\
\left.\frac{\partial f_{0}}{\partial r}\right|_{l,m+1/2}^{(k+1)}=\frac{f_{0,l+1/2,m+1/2}^{(k+1)}-f_{0,l-1/2,m+1/2}^{(k+1)}}{\Delta r_{l}}\\
\left.\frac{\partial f_{0}}{\partial\theta}\right|_{l+1,m+1/2}^{(k+1)}=\frac{f_{0,l+1,m+1}^{(k+1)}-f_{0,l+1,m}^{(k+1)}}{\Delta\theta_{m+1/2}}\\
\left.\frac{\partial f_{0}}{\partial\theta}\right|_{l,m+1/2}^{(k+1)}=\frac{f_{0,l,m+1}^{(k+1)}-f_{0,l,m}^{(k+1)}}{\Delta\theta_{m+1/2}}\end{array}\right.\]
 et \[
\left\{ \begin{array}{l}
\left.\frac{\partial f_{0}}{\partial r}\right|_{l+1/2,m+1}^{(k+1)}=\frac{f_{0,l+1,m+1}^{(k+1)}-f_{0,l,m+1}^{(k+1)}}{\Delta r_{l+1/2}}\\
\left.\frac{\partial f_{0}}{\partial r}\right|_{l+1/2,m}^{(k+1)}=\frac{f_{0,l+1,m}^{(k+1)}-f_{0,,m}^{(k+1)}}{\Delta r_{l+1/2}}\\
\left.\frac{\partial f_{0}}{\partial\theta}\right|_{l+1/2,m+1}^{(k+1)}=\frac{f_{0,l+1/2,m+3/2}^{(k+1)}-f_{0,l+1/2,m+1/2}^{(k+1)}}{\Delta\theta_{m+1}}\\
\left.\frac{\partial f_{0}}{\partial\theta}\right|_{i+1/2,m}^{(k+1)}=\frac{f_{0,l+1/2,m+1/2}^{(k+1)}-f_{0,l+1/2,m-1/2}^{(k+1)}}{\Delta\theta_{m}}\end{array}\right.\]
 ou $\Delta r_{l},$ $\Delta r_{l+1/2},$ $\Delta r_{l+1},$ $\Delta\theta_{m},$
$\Delta\theta_{m+1/2}$ et $\Delta\theta_{m+1}$ ont deja ete explicites
au paragraphe consacre a la definition des grilles.

En recombinant les differentes expressions on obtient

\begin{eqnarray*}
\left.S_{r}\right|_{l+1,m+1/2}^{(k+1)} & = & -D_{rr,l+1,m+1/2}\frac{f_{0,l+3/2,m+1/2}^{(k+1)}-f_{0,l+1/2,m+1/2}^{(k+1)}}{\Delta r_{l+1}}\\
 &  & -D_{r\theta,l+1,m+1/2}\frac{1}{r_{l+1}}\frac{f_{0,l+1,m+1}^{(k+1)}-f_{0,l+1,m}^{(k+1)}}{\Delta\theta_{m+1/2}}+F_{r,l+1,m+1/2}f_{0,l+1,m+1/2}^{(k+1)}\\
\left.S_{r}\right|_{l,m+1/2}^{(k+1)} & = & -D_{rr,l,m+1/2}\frac{f_{0,l+1/2,m+1/2}^{(k+1)}-f_{0,l-1/2,m+1/2}^{(k+1)}}{\Delta r_{l}}\\
 &  & -D_{r\theta,l,m+1/2}\frac{1}{r_{l}}\frac{f_{0,l,m+1}^{(k+1)}-f_{0,l,m}^{(k+1)}}{\Delta\theta_{m+1/2}}+F_{r,l,m+1/2}f_{0,l,m+1/2}^{(k+1)}\\
\left.S_{\theta}\right|_{l+1/2,m+1}^{(k+1)} & = & -D_{\theta r,l+1/2,m+1}\frac{f_{0,l+1,m+1}^{(k+1)}-f_{0,l,m+1}^{(k+1)}}{\Delta r_{l+1/2}}\\
 &  & -D_{\theta\theta,l+1/2,m+1}\frac{1}{r_{l+1/2}}\frac{f_{0,l+1/2,m+3/2}^{(k+1)}-f_{0,l+1/2,m+1/2}^{(k+1)}}{\Delta\theta_{m+1}}+F_{\theta,l+1/2,m+1}f_{0,l+1/2,m+1}^{(k+1)}\\
\left.S_{\theta}\right|_{l+1/2,m}^{(k+1)} & = & -D_{\theta r,l+1/2,m}\frac{f_{0,l+1,m}^{(k+1)}-f_{0,l,m}^{(k+1)}}{\Delta r_{l+1/2}}\\
 &  & -D_{\theta\theta,l+1/2,m}\frac{1}{r_{l+1/2}}\frac{f_{0,l+1/2,m+1/2}^{(k+1)}-f_{0,l+1/2,m-1/2}^{(k+1)}}{\Delta\theta_{m}}+F_{\theta,l+1/2,m}f_{0,l+1/2,m}^{(k+1)}\end{eqnarray*}


Comme la fonction est definie sur la demi-grille et les flux sur la
grille entiere, il faut egalement effectuer une interpolation, puisque
certaines valeurs de $f_{0}$ apparaissent sur la grille entiere dans
les relations precedentes. De maniere generale, sans se preoccuper
de l'expression precise du facteur de ponderation $\delta$ qui sera
discutee plus loin, on peut ecrire \[
\left\{ \begin{array}{l}
f_{0,l+1,m+1}^{(k+1)}=\left(1-\delta_{r,l+1,m+1}\right)f_{0,l+3/2,m+1}^{(k+1)}+\delta_{r,l+1,m+1}f_{0,l+1/2,m+1}^{(k+1)}\\
f_{0,l+1,m+1/2}^{(k+1)}=\left(1-\delta_{r,l+1,m+1/2}\right)f_{0,l+3/2,m+1/2}^{(k+1)}+\delta_{r,l+1,m+1/2}f_{0,l+1/2,m+1/2}^{(k+1)}\\
f_{0,l+1,m}^{(k+1)}=\left(1-\delta_{r,l+1,m}\right)f_{0,l+3/2,m}^{(k+1)}+\delta_{r,l+1,m}f_{0,l+1/2,m}^{(k+1)}\\
f_{0,l+1/2,m+1}^{(k+1)}=\left(1-\delta_{\theta,l+1/2,m+1}\right)f_{0,l+1/2,m+3/2}^{(k+1)}+\delta_{\theta,l+1/2,m+1}f_{0,l+1/2,m+1/2}^{(k+1)}\\
f_{0,l+1/2,m}^{(k+1)}=\left(1-\delta_{\theta,l+1/2,m}\right)f_{0,l+1/2,m+1/2}^{(k+1)}+\delta_{\theta,l+1/2,m}f_{0,l+1/2,m-1/2}^{(k+1)}\\
f_{0,l,m+1}^{(k+1)}=\left(1-\delta_{r,l,m+1}\right)f_{0,l+1/2,m+1}^{(k+1)}+\delta_{r,l,m+1}f_{0,l-1/2,m+1}^{(k+1)}\\
f_{0,l,m+1/2}^{(k+1)}=\left(1-\delta_{r,l,m+1/2}\right)f_{0,l+1/2,m+1/2}^{(k+1)}+\delta_{r,l,m+1/2}f_{0,l-1/2,m+1/2}^{(k+1)}\\
f_{0,l,m}^{(k+1)}=\left(1-\delta_{r,l,m}\right)f_{0,l+1/2,m}^{(k+1)}+\delta_{r,l,m}f_{0,l-1/2,m}^{(k+1)}\end{array}\right.\]


Comme on s'interesse principalement au cas de l'equation de Fokker-Planck
moyennee sur les trajectoires $\left(t\gg\tau_{t,b}\right)$, \[
\left\{ \frac{\partial\left(rp^{2}f_{0}\right)}{\partial t}\right\} +\left\{ p^{2}\frac{\partial\left(rS_{r}\right)}{\partial r}\right\} +\left\{ p^{2}\frac{\partial S_{\theta}}{\partial\theta}\right\} =0\]
 on elimine de fait la contribution polo\"{\i}dale en $\theta$ en
raison des moyennes effectuees:$\left\{ D_{r\theta}\right\} $ $=\left\{ D_{\theta r}\right\} =\left\{ D_{\theta\theta}\right\} =0$
et $\left\{ F_{\theta}\right\} $ $=0$ et on ne conserve que la dynamique
radiale. On pose par definition les moyennes des diffusion et convection
radiales sur les trajectoires des electrons $D_{rr}^{\left(0\right)}=\left\{ D_{rr}\right\} $
et $F_{r}^{\left(0\right)}=\left\{ F_{r}\right\} $ , ces valeurs
etant determinees en $B_{min}$. Dans cette limite, les equations
en jeu sont considerablement simplifiees car \[
\left\{ \frac{\partial\left(rp^{2}f_{0}\right)}{\partial t}\right\} +\left\{ p^{2}\frac{\partial\left(rS_{r}\right)}{\partial r}\right\} =0\]
 soit \[
\left\{ \frac{\partial\left(rp^{2}f_{0}\right)}{\partial t}\right\} +p^{2}\frac{\partial\left(r\left\{ S_{r}\right\} \right)}{\partial r}=\left\{ \frac{\partial\left(rp^{2}f_{0}\right)}{\partial t}\right\} +p^{2}\frac{\partial\left(rS_{r}^{\left(0\right)}\right)}{\partial r}=0\]
 avec \[
\mathbf{S}_{\mathbf{r}}^{(0)}=\left\{ \begin{array}{l}
S_{r}^{\left(0\right)}=\left\{ S_{r}\right\} \\
S_{\theta}^{\left(0\right)}=0\end{array}\right.\]
 On conserve ainsi toujours la forme conservative pour la composante
radiale \[
S_{\mathbf{r}}^{(0)}=-\overline{\overline{D}}_{r}^{(0)}\cdot\nabla_{r,\theta}f_{0}^{(0)}+\overline{\overline{F}}_{r}^{(0)}f_{0}^{(0)}\]
 avec les elements de la matrice de diffusion \[
D_{rr}^{(0)}=\left\{ D_{rr}\right\} \]
 \[
F_{r}^{(0)}=\left\{ F_{r}\right\} \]
 ou

\[
\left\{ S_{r}\right\} =S_{r}^{\left(0\right)}=\left\{ -D_{rr}\frac{\partial f}{\partial r}+F_{r}f\right\} =-\left\{ D_{rr}\right\} \frac{\partial f_{0}^{\left(0\right)}}{\partial r}+\left\{ F_{r}\right\} f_{0}^{\left(0\right)}=-D_{rr}^{\left(0\right)}\frac{\partial f_{0}^{\left(0\right)}}{\partial r}+F_{r}^{\left(0\right)}f_{0}^{\left(0\right)}\]
 Sous forme discrete, on obtient \[
r_{l+1/2}p_{i+1/2}^{2}\frac{f_{0,l+1/2}^{\left(0\right)(k+1)}-f_{0,l+1/2}^{\left(0\right)(k)}}{\Delta t}=-p_{i+1/2}^{2}\frac{r_{l+1}\left.S_{r}^{\left(0\right)}\right|_{l+1}^{(k+1)}-r_{l}\left.S_{r}^{\left(0\right)}\right|_{l}^{(k+1)}}{\Delta r_{l+1/2}}\]
 avec \begin{eqnarray*}
\left.S_{r}^{\left(0\right)}\right|_{l+1}^{(k+1)} & = & -D_{rr,l+1}^{\left(0\right)}\frac{f_{0,l+3/2}^{\left(0\right)(k+1)}-f_{0,l+1/2}^{\left(0\right)(k+1)}}{\Delta r_{l+1}}+F_{r,l+1}^{\left(0\right)}f_{0,l+1}^{\left(0\right)(k+1)}\\
\left.S_{r}^{\left(0\right)}\right|_{l}^{(k+1)} & = & -D_{rr,l}^{\left(0\right)}\frac{f_{0,l+1/2}^{\left(0\right)(k+1)}-f_{0,l-1/2}^{\left(0\right)(k+1)}}{\Delta r_{l}}+F_{r,l}^{\left(0\right)}f_{0,l}^{\left(0\right)(k+1)}\end{eqnarray*}
 et pour les interpolations entre les grilles \[
\left\{ \begin{array}{l}
f_{0,l+1}^{\left(0\right)(k+1)}=\left(1-\delta_{r,l+1}^{\left(0\right)}\right)f_{0,l+3/2}^{\left(0\right)(k+1)}+\delta_{r,l+1}^{\left(0\right)}f_{0,l+1/2}^{\left(0\right)(k+1)}\\
f_{0,l}^{\left(0\right)(k+1)}=\left(1-\delta_{r,l}^{\left(0\right)}\right)f_{0,l+1/2}^{\left(0\right)(k+1)}+\delta_{r,l}^{\left(0\right)}f_{0,l-1/2}^{\left(0\right)(k+1)}\end{array}\right.\]


Il est utile a ce stade d'exprimer sous forme matricielle le transport
radial, qui peut etre mis sous la forme generale

\[
r_{l+1/2}p_{i+1/2}^{2}\frac{f_{l+1/2}^{\left(0\right)(k+1)}-f_{l+1/2}^{\left(0\right)\left(k\right)}}{\Delta t}+p_{i+1/2}^{2}\sum_{l^{\prime}=l-1}^{l^{\prime}=l+1}\overline{\overline{M}}_{r,l^{\prime}+1/2}^{\left(0\right)}f_{0,l^{\prime}+1/2}^{\left(0\right)(k+1)}=0\]
 les coefficients de matrice $\overline{\overline{M}}_{r}^{\left(0\right)}$
valant,

\[
\overline{\overline{M}}_{r,l+3/2}^{\left(0\right)}=-\frac{r_{l+1}}{\Delta r_{l+1}\Delta r_{l+1/2}}D_{rr,l+1}^{\left(0\right)}+\frac{r_{l+1}}{\Delta r_{l+1/2}}F_{r,l+1}^{\left(0\right)}\left(1-\delta_{r,l+1}^{\left(0\right)}\right)\]
\begin{eqnarray*}
\overline{\overline{M}}_{r,l+1/2}^{\left(0\right)} & = & +\frac{r_{l+1}}{\Delta r_{l+1}\Delta r_{l+1/2}}D_{rr,l+1}^{\left(0\right)}+\frac{r_{l+1}}{\Delta r_{l+1/2}}F_{r,l+1}^{\left(0\right)}\delta_{r,l+1}^{\left(0\right)}\\
 &  & +\frac{r_{l}}{\Delta r_{l}\Delta r_{l+1/2}}D_{rr,l}^{\left(0\right)}-\frac{r_{l}}{\Delta r_{l+1/2}}F_{r,l}^{\left(0\right)}\left(1-\delta_{r,l}^{\left(0\right)}\right)\end{eqnarray*}
 \[
\overline{\overline{M}}_{r,l-1/2}^{\left(0\right)}=-\frac{r_{l}}{\Delta r_{l}\Delta r_{l+1/2}}D_{rr,l}^{\left(0\right)}-\frac{r_{l}}{\Delta r_{l+1/2}}F_{r,l}^{\left(0\right)}\delta_{r,l}^{\left(0\right)}\]


Il s'agit donc d'une simple matrice tridiagonale. Pour le choix de
$\delta_{r}^{\left(0\right)}$, la demarche est analogue a celle utilisee
pour $\delta_{\xi}$ puisque qu'il n'y a pas de contrainte particuliere
a prendre en compte sur la dependence radiale de la fonction de distribution
$f_{0}^{\left(0\right)}$. La fa\c{c}on la plus naturelle consiste
donc a effectuer une simple interpolation lineaire entre $f_{0,l+3/2}^{\left(0\right)(\infty)}$
et $f_{0,l+1/2}^{\left(0\right)(\infty)}$ ainsi qu'entre $f_{0,l+1/2}^{\left(0\right)(\infty)}$
et $f_{0,l-1/2}^{\left(0\right)(\infty)}$ selon les regles

\[
\left\{ \begin{array}{l}
f_{0,l+1}^{\left(0\right)(\infty)}=\left(\frac{\Delta r_{l+1/2}}{\Delta r_{l+1/2}+\Delta r_{l+3/2}}\right)f_{0,l+3/2}^{\left(0\right)(\infty)}+\left(\frac{\Delta r_{l+3/2}}{\Delta r_{l+1/2}+\Delta r_{l+3/2}}\right)f_{0,l+1/2}^{\left(0\right)(\infty)}\\
f_{0,l}^{\left(0\right)(\infty)}=\left(\frac{\Delta r_{l-1/2}}{\Delta r_{l-1/2}+\Delta r_{l+1/2}}\right)f_{0,l+1/2}^{\left(0\right)(\infty)}+\left(\frac{\Delta r_{l+1/2}}{\Delta r_{l-1/2}+\Delta r_{l+1/2}}\right)f_{0,l-1/2}^{\left(0\right)(\infty)}\end{array}\right.\]
 en tenant compte des pas de grille, ce qui donne \[
\left\{ \begin{array}{l}
\delta_{r,l+1}^{\left(0\right)}=\frac{\Delta r_{l+3/2}}{\Delta r_{l+1/2}+\Delta r_{l+3/2}}\\
\delta_{r,l}^{\left(0\right)}=\frac{\Delta r_{l+1/2}}{\Delta r_{l-1/2}+\Delta r_{l+1/2}}\end{array}\right.\]
 On retrouve bien que pour une grille uniforme, $\delta_{r,l+1}^{\left(0\right)}=\delta_{r,l}^{\left(0\right)}=1/2.$

Au centre du plasma en $r=0$, le flux $S_{r}^{\left(0\right)}$ doit
etre nul, puisque toutes les particules entrantes doivent ressortir.
Pour ce point singulier, on a donc \[
\left.S_{r}^{\left(0\right)}\right|_{l=0}^{(k+1)}=-D_{rr,l=0}^{\left(0\right)}\frac{f_{0,l=1/2}^{\left(0\right)(k+1)}-f_{0,l=-1/2}^{\left(0\right)(k+1)}}{\Delta r_{l}}+F_{r,l=0}^{\left(0\right)}f_{0,l=0}^{\left(0\right)(k+1)}=0\]
 et comme par symetrie $f_{0,l=1/2}^{\left(0\right)(k+1)}=f_{0,l=-1/2}^{\left(0\right)(k+1)},$
cela impose que \[
F_{r,l=0}^{\left(0\right)}=0\]
 car cette condition doit etre verifiee pour toute valeur de $f_{0,l=0}^{\left(0\right)(k+1)}$
qui n'est evidemment pas nulle au centre. Il n'y a evidemment aucune
contrainte sur la valeur de $D_{rr,l=0}^{\left(0\right)}$, car $\left.S_{r}^{\left(0\right)}\right|_{l=0}^{(k+1)}=0$
quelle que soit la valeur du coefficient de diffusion radiale.

Au bord en $r=a$, la condition usuelle est d'imposer que les particules
rapides quittent definitivement le plasma. La distribution est donc
la Maxwelienne non perturbee d'ou \[
f_{0,l=n_{r}}^{\left(0\right)(k+1)}=f_{M,l=n_{r}}^{\left(0\right)(k+1)}\]

